\documentclass[11pt]{article}
\usepackage[T1]{fontenc}
\usepackage[utf8]{inputenc}
\usepackage{lmodern,amsmath,amssymb,amsthm,booktabs,array}
\usepackage[margin=1in]{geometry}
\usepackage{microtype}
\usepackage[hidelinks]{hyperref}
\newtheorem{theorem}{Theorem}
\newtheorem{lemma}[theorem]{Lemma}
\newcommand{\N}{\mathbb N}
\newcommand{\R}{\mathbb R}
\newcommand{\Part}{\operatorname{Part}}
\setlength{\parindent}{0pt}
\setlength{\parskip}{0.5em}
\title{Disproof of Conjecture 00000000407\\[0.3em]
\large The stated equal-size Littlewood--Richardson table is zero}
\author{C0ldSmi1e}
\date{5 October 2026}
\begin{document}
\maketitle
\begin{abstract}
Both language versions index all three partitions by the same integer $n$.
For this exact indexing, every Littlewood--Richardson coefficient is zero
when $n>0$. We define and count conventional Littlewood--Richardson tableaux,
derive their size identity, and refute the complete positive-constant
asymptotic assertion. The accompanying Lean project proves these statements
for every $n$, including finiteness of the counted types. Independent exact
computations provide additional finite checks of the definitions.
\end{abstract}

\section{The exact assertion}
The official English statement is:
\begin{quote}
Conjecture: The number of nonzero entries of the LR coefficient table
$(c^\nu_{\lambda\mu})_{\lambda,\mu,\nu\vdash n}$ is asymptotically
$c\cdot4^n/n^{5/4}$ (a density estimate for LR-realizable triples among all
$4^n$ configurations).
\end{quote}
The Chinese text has the same simultaneous indexing and quantitative formula.
The full original bilingual file accompanies this submission as
\texttt{conjecture.md}. It is bound to upstream commit
\texttt{ba46ce0997cbf232fe95e5dced37ff4d32e1eb94} and source Git blob
\texttt{3439f91b71e05a0d57573e84e6c181d4c7a0ba86}.

Keep the indexing as printed, and write
\[
 A(n)=\#\{(\lambda,\mu,\nu)\in\Part(n)^3:
                   c^\nu_{\lambda\mu}\ne0\},\qquad
 g(n)=\frac{4^n}{n^{5/4}}\quad(n>0).
\]
The triples are ordered and each nonzero entry contributes one, regardless
of its coefficient value. The conventional asymptotic claim asserts that
some fixed real $c>0$ satisfies $A(n)/(c g(n))\to1$ as integers $n\to\infty$.

The parenthetical does not define an ambient configuration family. We do
not invent one: the explicit quantitative assertion is necessary whether
the parenthetical is explanatory or adds a further requirement. Refuting
that assertion therefore suffices. A source-only meaning assessment initially
deferred a proof of the entire statement because of this phrase, then cleared
this disproof-only scope. No claim is made about a different intended problem
obtained by changing the displayed indexing.

\section{Actual partitions and LR tableaux}
Represent an integer partition by its finite Young diagram, a lower set in
$\N^2$ for the coordinatewise order. Coordinates increase downward and to
the right. Diagram cardinality is partition size. Mathlib's
\texttt{YoungDiagram} has exactly this meaning; its row-length equivalence
identifies it with decreasing lists of positive parts. Thus the type
\texttt{Partition n} consists of all diagrams of cardinality $n$.

For diagrams $\lambda\subseteq\nu$, put $S=\nu\setminus\lambda$.
An LR tableau of shape $\nu/\lambda$ and content $\mu$ is a filling of $S$
with positive integers satisfying the following conditions:
\begin{enumerate}
\item Rows increase weakly from left to right; columns increase strictly downward.
\item Letter $i+1$ occurs exactly $\mu_{i+1}$ times, for every $i\in\N$,
      including all zero parts beyond the support of $\mu$.
\item Read rows from top to bottom, and each row from right to left.
      In every prefix, each positive letter occurs at least as often as its successor.
\end{enumerate}
The coefficient $c^\nu_{\lambda\mu}$ is the number of these tableaux;
if containment fails, that number is zero. This is the standard combinatorial
definition in \cite[Definition 2.1]{hahn} and \cite[Section 2]{gutschwager}.
The references identify the conventional objects; no external vanishing
theorem is assumed in the Lean proof.

The formal structure \texttt{LRTableau} records the filling, containment,
the two monotonicity conditions, all content equations, and the prefix
condition. Lean labels are shifted down by one: entry $i$ represents positive
letter $i+1$. No size identity is assumed as a field.

Pairwise row and column inequalities match the adjacent-cell definitions:
every skew row and column is an interval, and transitivity connects adjacent
comparisons. The formal reading relation is
\[
 (r,s)\preceq(r',s')\quad\Longleftrightarrow\quad
 r<r'\ \text{or}\ (r=r'\ \text{and}\ s'\le s).
\]
Its order properties are proved. The cells preceding an ending cell form
its inclusive reading prefix. Every nonempty prefix has a last cell; the
empty prefix satisfies the inequalities automatically. Consequently no
reading prefixes or labels are omitted.

\section{Finiteness and the size obstruction}
Every occurring Lean label $i$ has positive content, so
$i<\operatorname{height}(\mu)$. This is derived from the full content
equations. A tableau therefore injects into the finite function type
$S\to\{0,\ldots,\operatorname{height}(\mu)-1\}$. Its filling determines it;
the proof fields do not create extra objects. Thus the formal
\texttt{lrCoefficient}, defined using \texttt{Nat.card}, counts a finite
type and does not exploit the value of cardinality on infinite types.

Similarly, each coordinate of a cell in a size-$n$ lower set is less than
$n$: an occupied cell in row $i$ forces at least $i+1$ cells, and likewise
for columns. Every such diagram belongs to the powerset of an $n$ by $n$
grid. This proves finiteness of \texttt{Partition n}, its ordered triple
product, and the subtype of nonzero entries counted by \texttt{nonzeroCount}.

\begin{lemma}
Every LR tableau of shape $\nu/\lambda$ and content $\mu$ satisfies
$|\nu|=|\lambda|+|\mu|$.
\end{lemma}
\begin{proof}
Counting a diagram by row-coordinate fibers gives
$|\mu|=\sum_{i<\operatorname{height}(\mu)}\mu_{i+1}$.
Counting the skew cells by their entry values, using the derived label bound
and all content equations, gives
\[
 |\nu\setminus\lambda|
 =\sum_{i<\operatorname{height}(\mu)}
       \#\{a\in S:\operatorname{entry}(a)=i\}
 =\sum_{i<\operatorname{height}(\mu)}\mu_{i+1}=|\mu|.
\]
Containment and finite-set subtraction give
$|\nu\setminus\lambda|+|\lambda|=|\nu|$.
\end{proof}

\begin{theorem}
$A(n)=0$ for every positive integer $n$. No fixed positive real constant
$c$ satisfies $A(n)\sim c4^n/n^{5/4}$.
\end{theorem}
\begin{proof}
If an entry in the stated table were nonzero, there would be an LR tableau
with $|\lambda|=|\mu|=|\nu|=n$. The lemma would imply $n=n+n$, impossible
for $n>0$. The nonzero-entry type is empty, proving $A(n)=0$.

For completeness, the full ratio assertion is
\[
 \forall\varepsilon>0\ \exists N\in\N\ \forall n\ge N,
       \left|\frac{A(n)}{c g(n)}-1\right|<\varepsilon.
\]
Take $\varepsilon=1/2$ and, for any proposed $N$, take $n=N+1$.
Then $A(n)=0$ and $c g(n)>0$, making the absolute difference exactly $1$.
This contradicts $1<1/2$.
\end{proof}
In fact, $A(0)=1$: the empty tableau is the unique entry in the empty-partition
table. This finite initial value has no effect on the limit. The actual
normalized sequence $A(n)/g(n)$ is eventually zero and converges to zero.

\section{Formal statement and correspondence}
The real scale uses \texttt{Real.rpow} with $5/4$ explicitly in Lean's real
type, so there is no integer-division interpretation of the exponent. Its positivity for $n>0$ is proved.
\texttt{RatioAsymptotic c} is exactly the displayed epsilon--$N$ formula.
The final theorem is
\[
 \texttt{Conjecture407.conjecture407\_false}:\quad
 \neg\exists c\in\R,\ 0<c\ \land\ \texttt{RatioAsymptotic}(c).
\]
The stronger theorem \texttt{no\_real\_ratio\_constant} excludes every real
$c$ for the formal ratio predicate. Lean totalizes division at zero; that
additional case is immaterial to the ordinary positive-leading-constant
claim, whose denominator is positive on the tail.

\begin{center}\small
\begin{tabular}{@{}p{0.47\linewidth}p{0.46\linewidth}@{}}
\toprule
Formal declaration & Mathematical content\\
\midrule
\texttt{LRTableau.entry\_lt}, \texttt{lrTableauFinite}
 & Derived letter bound and finite tableau count\\
\texttt{diagram\_card\_eq\_sum\_rowLen}
 & Row-by-row diagram cardinality\\
\texttt{LRTableau.size\_balance}
 & Proved identity $|\nu|=|\lambda|+|\mu|$\\
\texttt{partitionFinite}, \texttt{nonzeroEntryFinite}
 & Finiteness of the exact index and counted-entry types\\
\texttt{lrCoefficient\_eq\_zero}, \texttt{nonzeroCount\_eq\_zero}
 & Every entry and its nonzero count vanish for $n>0$\\
\texttt{normalized\_tendsto\_zero}
 & Actual normalized limit in Mathlib's filter formulation\\
\texttt{conjecture407\_false}
 & Negation of the full positive-constant asymptotic claim\\
\bottomrule
\end{tabular}
\end{center}
\texttt{Audit407.lean} additionally proves the unit coefficient
$c^{\varnothing}_{\varnothing,\varnothing}=1$ and $A(0)=1$, $A(1)=A(2)=0$.
These are kernel-checked statements. No Schur-function or representation-theory
formalization is used or claimed; the Lean development uses the conventional
tableau definition directly.

\section{Independent finite checks and reproduction}
The separate standard-library Python verifier implements actual LR-tableau
enumeration with exact content checks, without a degree-mismatch shortcut.
It examines all $116\,947$ ordered triples for $0\le n\le10$, giving nonzero
counts $1,0,\ldots,0$. Its independently implemented sparse-polynomial
calculation uses Jacobi--Trudi \cite[(39)]{vanleeuwen} and alternant
coefficient extraction \cite[Bi-Alternant Corollary]{stembridge}.
Six variables suffice for the tested total input degree at most six.

All $4\,170$ tested triples agree between the two methods. Of these, $1\,110$
have compatible degree, with coefficient frequencies $816$ zeros, $293$ ones,
and one two; the other $3\,060$ have coefficient zero. Fourteen named
coefficient checks include units, Pieri cases, unequal input sizes and
$c^{(3,2,1)}_{(2,1),(2,1)}=2$. Five full-filling checks separately exercise
validity and isolated row, column, reading-word and content failures.
These finite experiments check the implementations; the proof for all $n$
is the Lean argument above and does not depend on them.

The project pins Lean 4.19.0 and Mathlib v4.19.0, including all nine exact
dependency revisions in \texttt{lean/lake-manifest.json}. From the submission
directory, reproduce the mathematical checks with:
\begin{verbatim}
cd lean
lake build
lake env lean -DwarningAsError=true Conjecture407.lean
lake env lean -DwarningAsError=true Audit407.lean
lake env lean -DwarningAsError=true Check.lean
cd ..
python3 auxiliary/verify.py > actual-output.json
cmp actual-output.json auxiliary/expected-output.json
\end{verbatim}
The verification directory supplies source identities, commands and outputs,
the full declaration/type/axiom inventory, and the independent review record.
Inspection tooling is separate from the mathematical modules. No admitted
proof, \texttt{native\_decide}, or custom axiom is used. Local validation and
the independent review are submission evidence; maintainer acceptance is a
separate decision.

\begin{thebibliography}{9}\small
\bibitem{hahn} H. Hahn, \emph{On Classical groups detected by the triple tensor
product and the Littlewood--Richardson semigroup}, Research in Number Theory
2, article 19 (2016), Definition 2.1.
\href{https://doi.org/10.1007/s40993-016-0049-3}{doi:10.1007/s40993-016-0049-3}.
\bibitem{gutschwager} C. Gutschwager, \emph{On Base Partitions and Cover
Partitions of Skew Characters}, Electronic Journal of Combinatorics 15
(2008), N30, Section 2.
\href{https://www.combinatorics.org/ojs/index.php/eljc/article/download/v15i1n30/pdf/}{Publisher PDF}.
\bibitem{vanleeuwen} M. A. A. van Leeuwen, \emph{Schur functions and alternating
sums}, Electronic Journal of Combinatorics 11(2) (2006), A5, equation (39).
\href{https://www.combinatorics.org/ojs/index.php/eljc/article/download/v11i2a5/pdf/}{Publisher PDF}.
\bibitem{stembridge} J. R. Stembridge, \emph{A Concise Proof of the
Littlewood--Richardson Rule}, Electronic Journal of Combinatorics 9 (2002),
N5, Bi-Alternant Corollary.
\href{https://www.combinatorics.org/ojs/index.php/eljc/article/download/v9i1n5/pdf}{Publisher PDF}.
\end{thebibliography}
\end{document}
